\documentclass[11pt]{article}
\usepackage{amsmath,amssymb,amsthm,booktabs,hyperref}
\newtheorem{theorem}{Theorem}
\newtheorem{lemma}[theorem]{Lemma}
\newtheorem{corollary}[theorem]{Corollary}
\title{Fourteen middle skeletons for bipartite quintic order32 graphs with an eight-geodesic}
\author{A conditional structural component of the F009 intake}
\date{}
\begin{document}\maketitle

\paragraph{Original problem, attribution and scope.}
The target belongs to Cambie, Goedgebeur, Jooken and Van den Eede,
\emph{On the order-diameter ratio of girth-diameter cages},
DMTCS28:4, SOFSEM2026\#1, arXiv:2511.21144v3,19August2026,
Appendix B, Table6, physical page23. Its parameter eight is diameter,
not an edge-cycle multiplicity. The independent-neighborhood/defect
viewpoint is borrowed from the manuscript and the retained F009 theory.
The star reduction was proposed by Root; the proof and integer normal
form below assemble standard bipartite, degree, omission-table and
distance arguments. No priority or novelty is asserted.

This theorem provides a finite normal form with fourteen \emph{middle
skeletons}, not fourteen ordinary graphs. It does not acquire any
current search family, certify the reference count of $16\,467$
bipartite positive classes, or prove the full $25\,408$-class census.
No Source, graph/state/gzip body, generator, control or TeX compiler
was executed. Ordinary isomorphism, including root and color changes,
must still be identified as stated at the end.

\paragraph{Ambient hypotheses.}
Let $G$ be a connected simple bipartite $5$-regular graph on $32$
vertices, with a fixed geodesic $p_0,\ldots,p_8$ of length eight.
We do not initially assume a global diameter upper bound.
The endpoints and $p_4$ lie in the same color class, called $U$;
write $V$ for the other class. Define
\[
 A=N(p_0),\ B=N(p_4),\ C=N(p_8),\ I=A\cup B\cup C,
 \quad J=V(G)\setminus I,\quad F=G[J].
\]
For $K=J\setminus\{p_0,p_4,p_8\}$ put
\[
 L=\{x\in K:N(x)\cap A\ne\varnothing\},\quad
 R=\{x\in K:N(x)\cap C\ne\varnothing\},\quad
 M=K\setminus(L\cup R).
\]
The complete graph, not a partial record's absent edges, defines these sets.

\begin{lemma}[Forced star and middle data]\label{star}
Both color classes have size $16$. There is a unique $z\in V\setminus I$,
and $F$ is the star $K_{1,5}$ centered at $z$. The three centers
$p_0,p_4,p_8$ are isolated in $F$, and $z\in M$.
There are integers $t,u\geq0$ such that
\begin{itemize}
\item $t$ vertices of $M\setminus\{z\}$ are isolated in $F$ and have
neighborhood $B$;
\item $u$ vertices of $M\setminus\{z\}$ meet $z$, and each meets
exactly four vertices of $B$, missing exactly one;
\item the other $5-u$ neighbors of $z$ all lie in one side, which
we name $L$.
\end{itemize}
Put $r=t+u$ and $q=|M|=1+r$. Then
\[
 |L|,|R|\geq5,\qquad |L|+|R|=13-r,\qquad r\leq3.   \tag{1}
\]
Every $b_i\in B$ has a nonempty neighbor set in exactly one side.
If $m_i$ is the number of the $u$ nonisolated middle vertices missing
$b_i$, then its side arity is
\[
 s_i=4-r+m_i,\qquad \sum_{i=1}^5m_i=u,\qquad1\leq s_i\leq4. \tag{2}
\]
Both side-port parts $B_L,B_R$ are nonempty. Moreover $u>0$.
\end{lemma}
\begin{proof}
Regularity makes the two color classes equal. The geodesic center
distances $4,4,8$ make the neighborhoods disjoint, so $I$ is a
$15$-subset of $V$. Thus $J=U\sqcup\{z\}$, and its only possible
internal edges are the five edges at $z$. No center meets $z$, or
it would have $z$ in its already specified neighborhood. Since
$z$ is in the same color as $A,C$, it has no neighbor in either
and belongs to $M$.

An internal-star leaf in $M$ has four remaining neighbors, all in
$B$. A non-leaf in $M\setminus\{z\}$ has all five neighbors in
$B$. A neighbor of $z$ cannot lie in both sides: a leaf in each
would give a marked-root path of length six. Similarly no $B$-vertex
meets both sides. Every endpoint-neighborhood row has four side
neighbors, so each side has at least four vertices; equality would
close a proper $K_{5,5}$ component. Hence both have at least five.
Counting the $14$ noncenter vertices of $J$ gives (1). In particular
$u\leq r\leq3$, so $5-u\geq2$: a side containing the remaining
star leaves exists and fixes the naming $L$.

The core vertices adjacent to $b_i$ are $p_4$, the $t$ isolated
middle vertices and $u-m_i$ other middle vertices. Degree five
then gives (2); its bounds use $m_i\leq u\leq r\leq3$.
The side part $B_R$ cannot be empty, since $R,C,p_8$ would be a
proper component. Put $n_L=|L|,n_R=|R|$. Degree sums on the sides give
\[
 \sum_{B_L}s_i=5n_L-25+u,\qquad
 \sum_{B_R}s_i=5n_R-20.                             \tag{3}
\]
For the first equation there are twenty $A$ incidences and $5-u$
star incidences; for the second there are twenty $C$ incidences
and no star incidences. If $u>0$, the first sum is positive by
$n_L\geq5$. If $u=0$ and $B_L$ were empty, the whole left side,
$A,p_0,z$ would be a proper component. Thus both port parts are
nonempty in every case.

Finally suppose $u=0$. All side arities equal $s=4-t\in\{1,2,3,4\}$.
The right equation in (3) gives
\[
 s|B_R|=5(n_R-4).
\]
Coprimality of $s$ and $5$ forces $5\mid |B_R|$, contrary to
$1\leq |B_R|\leq4$. Thus $u>0$.
\end{proof}

\begin{theorem}[The fourteen skeletons]\label{fourteen}
Up to permuting the vertices of $B$ and the middle vertices with
the same role, the complete possibilities for the middle and
side-size data are the fourteen rows of Table~\ref{table}.
The tuples $m_L,m_R$ list all missing multiplicities in the indicated
port part, including zeros; their lengths give the numbers of ports.
The side $L$ is distinguished by the $5-u$ star leaves. These rows
are not ordinary graph classes or multiplicities of such classes.
\end{theorem}
\begin{table}[h]\centering\small
\begin{tabular}{c|ccc|cc|l|l}
\toprule
row&$q$&$t$&$u$&$n_L$&$n_R$&$m_L$&$m_R$\\\midrule
1&2&0&1&6&6&$(0,0)$&$(1,0,0)$\\
2&3&1&1&6&5&$(0,0,0)$&$(1,0)$\\
3&3&0&2&5&6&$(0)$&$(2,0,0,0)$\\
4&3&0&2&5&6&$(0)$&$(1,1,0,0)$\\
5&3&0&2&6&5&$(1,0,0)$&$(1,0)$\\
6&4&2&1&5&5&$(0)$&$(1,0,0,0)$\\
7&4&1&2&5&5&$(1)$&$(1,0,0,0)$\\
8&4&1&2&5&5&$(0,0)$&$(2,0,0)$\\
9&4&1&2&5&5&$(0,0)$&$(1,1,0)$\\
10&4&0&3&5&5&$(2)$&$(1,0,0,0)$\\
11&4&0&3&5&5&$(1,0)$&$(2,0,0)$\\
12&4&0&3&5&5&$(1,0)$&$(1,1,0)$\\
13&4&0&3&5&5&$(0,0,0)$&$(3,0)$\\
14&4&0&3&5&5&$(0,0,0)$&$(2,1)$\\\bottomrule
\end{tabular}
\caption{All middle/port skeletons before endpoint incidences are chosen.}
\label{table}\end{table}
\begin{proof}
Write $b=|B_L|$ and $k=\sum_{B_L}m_i$. Equation (3) is
\[
 (4-r)b+k=5(n_L-5)+u,\quad 0\leq k\leq u,\quad
 n_L+n_R=13-r.                                    \tag{4}
\]
The right sum in (3) is at least five whereas each port arity is
at most four, so $|B_R|\geq2$ and $1\leq b\leq3$.
For $r=1$, necessarily $u=1$: (4) gives only $n_L=n_R=6,b=2,k=0$.
For $r=2,u=1$ it gives only $(n_L,n_R,b,k)=(6,5,3,0)$.
For $r=2,u=2$ it gives $(5,6,1,0)$ or $(6,5,3,1)$.
For $r=3$ both sides have size five, and $b+k=u$.
Thus for $u=1$ only $(b,k)=(1,0)$ occurs; for $u=2$ these
are $(1,1),(2,0)$; for $u=3$ these are $(1,2),(2,1),(3,0)$.

Distribute $k$ missing middle vertices among $b$ ports and $u-k$
among $5-b$ ports, up to permutation inside each part. For a total
two the partitions are $(2)$ or $(1,1)$ when enough positions are
available; for a total three on two positions they are $(3,0)$
or $(2,1)$. The displayed cases give respectively $1,1,3,1,3,5$
rows for $(q,u)=(2,1),(3,1),(3,2),(4,1),(4,2),(4,3)$.
There is no further middle incidence datum: each of the $u$ middle
leaves misses exactly one port. Its missing port determines its
entire core neighborhood. This proves the table and its equivalence
under the stated role-preserving permutations.
\end{proof}

\paragraph{The explicit integer normal form.}
Select one row of Table~\ref{table}. Take disjoint labeled sets
$L,R$ of the displayed sizes, port sets $B_L,B_R$ of the displayed
tuple lengths, endpoint-neighborhood sets $A,C$ of size five,
and the following additional vertices:
\[
 p_0,p_8,z,\qquad X_0=\{p_4,h_1,\ldots,h_t\},\qquad
 X_i\ (i\in B),\quad |X_i|=m_i.
\]
The union $X=X_0\sqcup\bigsqcup_iX_i$ is the central core.
Every $X_0$ vertex meets all five ports. Every vertex of $X_i$
meets $z$ and all ports other than $b_i$. Choose a subset
$S\subset L$ of size $5-u$, and join $z$ to $S$. These are all
the neighbors of $z$.

For each port $b_i$, choose a \emph{set} $T_i$ of $s_i=4-r+m_i$
distinct vertices of its assigned side, and add exactly those side
edges. Equivalently, choose a zero--one port-by-side incidence table
with those exact row sums. Define on the two sides
\[
 h_L(v)=|\{i\in B_L:v\in T_i\}|+\mathbf1_{v\in S},\qquad
 h_R(v)=|\{i\in B_R:v\in T_i\}|.                   \tag{5}
\]
We impose $h_D(v)\leq4$ at every side vertex, to retain a positive
endpoint-neighborhood degree. This bound is in fact automatic in
the table: $|B_L|\leq3$ and $|B_R|\leq4$.

Each of the five endpoint rows has four distinct neighbors in its
side. Describe its omitted vertices as follows:
\begin{itemize}
\item On a side of size five, omit singleton $v$ exactly $h_D(v)$
times. Thus the five omission rows are fixed by their column counts,
up to permutation of the endpoint-neighborhood vertices.
\item On a side of size six, choose any loopless multigraph with
five edges and degree sequence $h_D(v)$. Its five edges, with
parallel copies retained, are the omitted pairs of the five rows.
\end{itemize}
Join $p_0$ to all of $A$ and $p_8$ to all of $C$. Each endpoint row
meets precisely the complement of its omission row in its side.
There are no edges beyond the listed ones. These instructions,
(5), and the table give the entire graph without a large graph search.

\begin{theorem}[Necessity, realization and converse]\label{converse}
Every graph in the ambient hypotheses has the preceding normal form.
Conversely every such choice constructs a connected simple bipartite
$5$-regular graph on $32$ vertices of girth exactly four and diameter
exactly eight, with the displayed $p_4$ at the middle of a diametral
geodesic from $p_0$ to $p_8$. Every row of Table~\ref{table} admits
at least one such realization. The converse does not require a later
diameter-filter computation.
\end{theorem}
\begin{proof}[Incidence necessity and existence]
In an actual graph, the middle incidences are precisely those proved
in Lemma~\ref{star}. All other vertices in $U$ are $p_0,p_8,L,R$;
they have no mutual edges by bipartiteness. At a side vertex,
the degree available for endpoint rows is $5-h_D(v)$. Therefore
its number of omissions among the five rows is exactly $h_D(v)$.
Each row omits one or two vertices according to side size five or
six, proving the asserted table/multigraph description. Conversely
these column degrees give exactly the required degree at every side
vertex and a positive endpoint-neighborhood incidence.

The balances (3) imply $\sum_v h_D(v)=5(n_D-4)$, which is five
or ten. The singleton omission rows therefore exist when $n_D=5$.
When $n_D=6$ their degree sequence has total ten and maximum at
most four. A nonnegative integer sequence of total $2m$ is the
degree sequence of a loopless multigraph if its maximum is at most
$m$: arrange equal-vertex stubs in consecutive blocks on a circle
of $2m$ positions and pair opposite positions. No block of length
at most $m$ contains two opposite positions, so no loop occurs.
This gives the five required omission pairs for $m=5$.

For every skeleton, each prescribed port arity is between one and
four, so distinct-vertex subsets of the assigned side exist. Any
choice and any $S$ of the prescribed size has $h_D\leq4$ by the
port-count bounds above. Hence the omission constructions apply
and every skeleton is realized; this uses no enumeration.

Every endpoint row has its center and four side neighbors. Each
port has $1+t+u-m_i$ core neighbors and $4-r+m_i$ side neighbors,
total five. The core and $z$ have degree five by construction,
and so do the side vertices by (5). The two color classes are
\[
 U=\{p_0,p_8\}\sqcup L\sqcup R\sqcup X,
 \qquad V=A\sqcup B\sqcup C\sqcup\{z\},
\]
each of size sixteen, and all edges go between them. Simplicity
uses sets for the port incidences and distinct rows for repeated
omission pairs. Repeated rows do not create multiple graph edges.

The core is connected through $B$, since $X_0$ is nonempty and
each nonisolated core vertex has four port neighbors. The vertex
$z$ meets $u>0$ core vertices. Every side vertex meets its own
endpoint neighborhood, every port has a side neighbor, and both
port parts are nonempty. Thus the whole graph is connected.
Any vertex in $U$ has five neighbors in $V$; their twenty other
$U$-incidences cannot all be distinct among only fifteen other
$U$ vertices. Two share a second neighbor in $U$, giving a
four-cycle. Thus the bipartite graph has girth exactly four.
\end{proof}

\begin{proof}[Diameter and the geodesic]
Any two core vertices have a common port neighbor: each sees at
least four of the five ports. Distinct core vertices are therefore
at distance two.

An endpoint center is at distance exactly four from \emph{every}
core vertex. For $p_0$ and a vertex of $X_0$, use any left port
and one of its side neighbors, which meets an endpoint row. For
$p_0$ and a nonisolated core vertex, use
\[
 p_0,A,S,z,X_i.
\]
For $p_8$ and $X_0$, use a right port. For $p_8$ and a vertex of
$X_i$, choose a right port other than its one missing port; this
is possible because $|B_R|\geq2$. These give length-four paths.
No endpoint and core vertex has a common neighbor, since endpoint
neighborhoods are $A,C$, whereas core neighborhoods are in
$B\cup\{z\}$. Their common color makes distance four minimal.
Every endpoint false twin has the same distances to the core.

Every $A$-row meets a star-side leaf. Indeed, using $n_R\geq5$,
\[
 |L\setminus S|=n_L-(5-u)\leq3-t<4,
\]
whereas an $A$-row uses four side vertices. It consequently has
distance three to some nonisolated core vertex, through $S,z$.
Every $C$-row meets a side vertex with a right-port neighbor.
Otherwise the union $T$ of all right-port supports would lie in
its $n_R-4$ omitted vertices. Their combined port degree is at
most $4(n_R-4)$, contrary to the required total $5(n_R-4)$ from
(3). Thus a $C$-row is within three of $X_0$. Ports and $z$ are
within one of a core vertex. Hence every vertex of $V$ is within
three of some core vertex.

A side vertex with a port neighbor is within two of $X_0$.
A star-side leaf is within two of a nonisolated core vertex,
through $z$. Any remaining side vertex has neither a port nor a
star neighbor. It has all five endpoint-row neighbors and is a
false twin of its endpoint center. Therefore every vertex of $U$
is either a core vertex, an endpoint false twin at distance four
from \emph{all} core vertices, or is within two of some core vertex.
The endpoint-false-twin exception is necessary; these vertices
need not be within two of the core.

These bounds prove all-pairs distance at most eight. Two endpoint
twins can use the same core vertex, at cost $4+4$; an endpoint
twin and any other vertex cost at most $4+3$ or $4+2$. Two
$V$ vertices cost at most $3+2+3$ using their possibly different
core vertices, and the remaining pairs cost no more.

Finally, deleting $X$ separates $p_0,L,A,z,B_L$ from $p_8,R,C,B_R$.
Every endpoint is distance four from $X$, so the marked endpoint
distance is at least eight. It is at most eight through $p_4\in X_0$,
and both halves have length four. Hence $p_0,\ldots,p_8$ can be
chosen with that middle, and the global diameter is exactly eight.
The recovered middle set is precisely $\{z\}\sqcup(X\setminus\{p_4\})$,
with the claimed $t,u$ roles, because every other side vertex has
a positive endpoint-neighborhood incidence. This completes the converse.
\end{proof}

\paragraph{Ordinary isomorphism and the remaining count.}
For a complete graph identify vertices with equal open neighborhoods,
give each class its cardinality, and retain class adjacency. Distinct
twins are nonadjacent, and class adjacency is complete or empty, so
independent expansion recovers the graph. An ordinary graph isomorphism
induces a weight-preserving isomorphism of this full quotient, and
conversely such a quotient isomorphism lifts. This intrinsic quotient
is therefore a complete ordinary-isomorphism test.

Some visible classes are $X_0$ of weight $t+1$, each nonempty $X_i$
of weight $m_i$, repeated endpoint omission rows, repeated identical
port rows, and side vertices with the same complete incidences.
These suggestions are not a replacement for the \emph{full} quotient.
Different middle rows, root choices, endpoint orientations, twin
labels and even color-class orientations can describe the same
ordinary graph. No division by a twin factorial, a port-permutation
order, fourteen, or a guessed free action is justified.

The result replaces arbitrary graph completion in this bipartite
stratum by zero--one port tables and five singleton/pair omission
rows on sides of size at most six. A complete ordinary class count
would still require an exhaustive independent enumeration of these
finite integer forms and exact identification of their intrinsic
weighted quotients, or an equivalent orbit proof. Neither was
performed here. The manuscript's reported counts are comparisons,
not input stopping rules or proof conditions.

\paragraph{A general-degree necessity statement, not a new classification.}
For a connected simple bipartite $k$-regular graph on $6k+2$ vertices
with a fixed eight-geodesic, $k\geq3$, the same setup gives color
classes $3k+1$, $I$ of size $3k$, and an extra vertex $z$ centering
$F=K_{1,k}$. The middle consists of $z$, $t$ isolated vertices and
$u$ star leaves missing one port each. Necessarily
\[
 q=1+t+u\leq k-1,\quad n_L,n_R\geq k,\quad
 n_L+n_R=3k-2-t-u,
\]
and, after naming the star side $L$,
\[
 m_i\geq0,\quad\sum m_i=u,\quad s_i=k-q+m_i,
\quad \sum_{B_L}s_i=k(n_L-k)+u,
\quad \sum_{B_R}s_i=k(n_R-k+1).                    \tag{6}
\]
These are proved by exactly the color balance, proper-component,
star and degree-sum arguments above. Both port parts are nonempty:
if $u>0$ use the left positive balance; if $u=0$ use connectedness.
When $u=0$, put $s=k-1-t\in\{1,\ldots,k-1\}$. If
$\gcd(k,s)=1$, the right balance forces $k\mid|B_R|$, impossible.
In particular prime degree rules out $u=0$.

For composite $k$, coprimality need not hold. The arithmetic pattern
$k=4,t=1,u=0$, side arity two, $|B_R|=2,n_R=4,n_L=5$
satisfies these necessary balances. This is not an assertion of a
realized graph or diameter-eight converse. Endpoint omissions in
general have size $n_D-k+1$, which need not be singleton or pair
rows. No general-degree template classification, count or converse
is claimed; those require their own hypotheses and proof.
\end{document}
